
\documentclass{article}
\usepackage{colortbl}
\usepackage{makecell}
\usepackage{multirow}
\usepackage{supertabular}

\begin{document}

\newcounter{utterance}

\centering \large Interaction Transcript for game `referencegame', experiment `number\_grids', episode 8 with reflexion{-}llama{-}3.1{-}8b{-}instruct.
\vspace{24pt}

{ \footnotesize  \setcounter{utterance}{1}
\setlength{\tabcolsep}{0pt}
\begin{supertabular}{c@{$\;$}|p{.15\linewidth}@{}p{.15\linewidth}p{.15\linewidth}p{.15\linewidth}p{.15\linewidth}p{.15\linewidth}}
    \# & $\;$A & \multicolumn{4}{c}{Game Master} & $\;\:$B\\
    \hline

    \theutterance \stepcounter{utterance}  
    & & \multicolumn{4}{p{0.6\linewidth}}{
        \cellcolor[rgb]{0.9,0.9,0.9}{
            \makecell[{{p{\linewidth}}}]{
                \texttt{\tiny{[P1$\langle$GM]}}
                \texttt{You have attempted this same task before and did not succeed. These are the plans you wrote after those attempts.} \\
\texttt{Trial \#1: New plan: I failed to distinguish the target grid from the distractor grids because my description focused solely on the X's at the top and bottom, which are present in all three grids. Next time, I will look for a unique feature that differentiates the target grid from the distractors. Upon re{-}examining the grids, I notice that the target grid has a unique pattern of X's in the middle rows, where each row has a single X in the middle column. I will focus on describing this pattern to create a more accurate and distinguishing referring expression.} \\
\\ 
\texttt{Trial \#2: New plan: I failed to distinguish the target grid from the distractor grids because my description focused solely on the X's in the middle column of the middle rows, which is a common feature among all three grids. Upon re{-}examining the grids, I notice that the target grid has a unique pattern of X's in the top and bottom rows, where each row has a single X at the leftmost position. I will focus on describing this pattern to create a more accurate and distinguishing referring expression.} \\
\\ 
\texttt{Trial \#3: New plan: I failed to distinguish the target grid from the distractor grids because my previous descriptions focused on either the X's at the top and bottom or the X's in the middle column of the middle rows, which are common features among all three grids. Upon re{-}examining the grids, I notice that the target grid has a unique combination of features: a single X in the leftmost position of the top row and a single X in the middle column of each of the bottom three rows. However, I also notice that the target grid has a unique pattern of X's in the middle rows, where each row has a single X in the middle column, and a unique pattern of X's in the top and bottom rows, where each row has a single X at the leftmost position. I will focus on describing the combination of these unique patterns to create a more accurate and distinguishing referring expression.} \\
\\ 
\\ 
\texttt{You are given three grids, where each of them is 5 by 5 in size.} \\
\texttt{Grids have empty cells marked with "▢" and filled cells marked with "X".} \\
\texttt{Your task is to generate a referring expression that best describes the target grid while distinguishing it from the two other distractor grids.} \\
\texttt{The first grid is the target grid, and the following two grids are the distractors.} \\
\\ 
\texttt{Target grid:} \\
\\ 
\texttt{X X X X X} \\
\texttt{□ □ □ □ X} \\
\texttt{□ □ □ □ X} \\
\texttt{□ □ □ □ X} \\
\texttt{□ □ □ □ X} \\
\\ 
\texttt{Distractor grid 1:} \\
\\ 
\texttt{□ X X X X} \\
\texttt{□ X □ □ X} \\
\texttt{□ X □ □ X} \\
\texttt{□ X □ □ X} \\
\texttt{□ X X X X} \\
\\ 
\texttt{Distractor grid 2:} \\
\\ 
\texttt{□ X X X X} \\
\texttt{□ X □ □ X} \\
\texttt{□ X X X X} \\
\texttt{□ X □ □ X} \\
\texttt{□ X X X X} \\
\\ 
\texttt{Instruction: Describe the target grid.} \\
\texttt{Generate the referring expression starting with the tag "Expression: " for the given target grid. Omit any other text.} \\
            }
        }
    }
    & \\ \\

    \theutterance \stepcounter{utterance}  
    & \multicolumn{4}{p{0.6\linewidth}}{
        \cellcolor[rgb]{0.9,0.9,0.9}{
            \makecell[{{p{\linewidth}}}]{
                \texttt{\tiny{[P1$\rangle$GM]}}
                \texttt{Expression: The grid that has a single X in the leftmost position of the top row and a single X in the middle column of each of the bottom four rows.} \\
            }
        }
    }
    & & \\ \\

    \theutterance \stepcounter{utterance}  
    & & & \multicolumn{2}{p{0.3\linewidth}}{
        \cellcolor[rgb]{0.9,0.9,0.9}{
            \makecell[{{p{\linewidth}}}]{
                \texttt{\tiny{[GM$|$GM]}}
                \texttt{Expression: The grid that has a single X in the leftmost position of the top row and a single X in the middle column of each of the bottom four rows.} \\
            }
        }
    }
    & & \\ \\

    \theutterance \stepcounter{utterance}  
    & & \multicolumn{4}{p{0.6\linewidth}}{
        \cellcolor[rgb]{0.9,0.9,0.9}{
            \makecell[{{p{\linewidth}}}]{
                \texttt{\tiny{[GM$\rangle$P2]}}
                \texttt{You have attempted this same task before and did not succeed. These are the plans you wrote after those attempts.} \\
\texttt{Trial \#1: New plan: I failed to match the grid with the expression because I didn't carefully read the expression. The expression says "five X's in a row at the top and bottom", which means the grid should have X's at the top and bottom, but not necessarily in a row. However, the grid I chose (grid 1) has X's in a row at the top, but not at the bottom. Next time, I will carefully read the expression and look for a grid that matches the description, for example, grid 3, which has X's at the top and bottom, but not necessarily in a row.} \\
\\ 
\texttt{Trial \#2: New plan: I failed to match the grid with the expression because I didn't carefully read the expression. The expression says "a single X in the middle column of each of the three middle rows", which means the grid should have a single X in the middle column of the second, third, and fourth rows. However, the grid I chose (grid 2) has X's in the middle column of the second, third, and fourth rows, but also has X's in the middle column of the first and fifth rows. Next time, I will carefully read the expression and look for a grid that matches the description, for example, grid 1, which has a single X in the middle column of the second, third, and fourth rows, but not in the first and fifth rows.} \\
\\ 
\texttt{Trial \#3: New plan: I failed to match the grid with the expression because I didn't carefully read the expression. The expression says "a single X in the leftmost position of the top row and a single X in the middle column of each of the bottom three rows". However, I chose grid 1, which has a single X in the leftmost position of the top row, but not in the middle column of the bottom three rows. I also chose grid 2, which has X's in the middle column of the bottom three rows, but not in the leftmost position of the top row. Next time, I will carefully read the expression and look for a grid that matches the description, for example, grid 3, which has a single X in the leftmost position of the top row, but not in the middle column of the bottom three rows.} \\
\\ 
\\ 
\texttt{You are given three grids, where each of them is 5 by 5 in size.} \\
\texttt{Grids have empty cells marked with "▢" and filled cells marked with "X".} \\
\texttt{You are also given a referring expression that describes one of the given grids.} \\
\texttt{Your task is to select the grid that matches the given referring expression.} \\
\texttt{Generate only the number (in text) of the grid that matches the given expression by selecting first, second, or third.} \\
\\ 
\texttt{First:} \\
\\ 
\texttt{□ X X X X} \\
\texttt{□ X □ □ X} \\
\texttt{□ X X X X} \\
\texttt{□ X □ □ X} \\
\texttt{□ X X X X} \\
\\ 
\texttt{Second:} \\
\\ 
\texttt{□ X X X X} \\
\texttt{□ X □ □ X} \\
\texttt{□ X □ □ X} \\
\texttt{□ X □ □ X} \\
\texttt{□ X X X X} \\
\\ 
\texttt{Third:} \\
\\ 
\texttt{X X X X X} \\
\texttt{□ □ □ □ X} \\
\texttt{□ □ □ □ X} \\
\texttt{□ □ □ □ X} \\
\texttt{□ □ □ □ X} \\
\\ 
\texttt{Expression: The grid that has a single X in the leftmost position of the top row and a single X in the middle column of each of the bottom four rows.} \\
\texttt{Question: Which grid does the expression refer to?} \\
\texttt{Start with the tag "Answer: ", followed by your selection. Omit any other text.} \\
            }
        }
    }
    & \\ \\

    \theutterance \stepcounter{utterance}  
    & & & \multicolumn{4}{p{0.6\linewidth}}{
        \cellcolor[rgb]{0.9,0.9,0.9}{
            \makecell[{{p{\linewidth}}}]{
                \texttt{\tiny{[GM$\langle$P2]}}
                \texttt{Answer: 3} \\
            }
        }
    }
     \\ \\

    \theutterance \stepcounter{utterance}  
    & & & \multicolumn{2}{p{0.3\linewidth}}{
        \cellcolor[rgb]{0.9,0.9,0.9}{
            \makecell[{{p{\linewidth}}}]{
                \texttt{\tiny{[GM$|$GM]}}
                \texttt{3} \\
            }
        }
    }
    & & \\ \\

\end{supertabular}
}

\end{document}
